\section{Study 3: PyBullet + Elevated Combined}
\label{sec:study3}

\subsection{Motivation}

Studies 1 and 2 establish complementary capabilities in isolation.
Study 1 (\texttt{mc-pilot-elevated/}) demonstrates 100\% hit rates across
$z \in \{0.5, 1.0, 1.5, 2.0\}$\,m using the NumPy ballistic simulator.
Study 2 (\texttt{mc-pilot-pybullet/}) integrates a KUKA iiwa7 arm in PyBullet with
velocity-slip noise.
Neither folder provides what an end-to-end demonstration requires: a physically
rendered arm throw from a non-standard release height.
Study 3 fills this gap, combining the elevated parameter sets from Study 1 with the
PyBullet physics pipeline from Study 2.
The sub-project \texttt{mc-pilot-pb-elevated/} introduces no new algorithm
experiments; its purpose is visualisable validation of the geometry-corrected
parameters under a higher-fidelity simulator.

\subsection{Engineering Challenges}
\label{sec:study3:challenges}

Two mechanical problems arise when lifting the release point above ground.

\paragraph{Joint velocity cap.}
The KUKA iiwa7 hardware joint velocity limits are
$[1.71,\, 1.71,\, 1.75,\, 2.27,\, 2.44,\, 3.14,\, 3.14]$\,rad/s.
For multi-joint coordinated motion, these limits cap the achievable end-effector
(EE) linear speed at roughly 2.5\,m/s.
The elevated configs from Study 1 require $u_M = 3.5$\,m/s.
The apparent conflict is harmless because the ball velocity at release is set
explicitly via PyBullet's \texttt{resetBaseVelocity}, independently of the arm
kinematics.
The arm trajectory is cosmetic: it demonstrates the physical motion, but any
discrepancy between commanded EE velocity and actual joint tracking has zero effect
on the ball's initial state.
The cap is enforced inside \texttt{ArmController.plan\_throw()} through a Jacobian
pseudoinverse scaling step.
Relaxing this check via \texttt{vel\_limit\_multiplier=1.5} prevents spurious
warnings without altering ball physics.

\paragraph{EE reach at elevated heights.}
The iiwa7 EE at its neutral pose sits approximately 0.71\,m above the arm base.
Release heights of 1.0, 1.5, and 2.0\,m place the IK target 0.29--1.29\,m above
the natural EE position when the arm is mounted on the ground plane.
This exceeds the arm's reach.
The resolution is to mount the arm on a pedestal: setting \texttt{base\_position}
$= (0, 0, z_{\mathrm{release}} - 0.5)$ places the neutral EE height at
$z_{\mathrm{release}} + 0.21$\,m, so the IK target at $z_{\mathrm{release}}$ lies
approximately 0.21\,m below the neutral configuration---well within the arm's
kinematic range.
Table~\ref{tab:pbe_configs} summarises the per-height parameter choices.

\begin{table}[h]
  \centering
  \caption{Per-configuration parameters for Study 3 (PBE-B/C/D).
           All three share $u_M = 3.5$\,m/s, $N_{\mathrm{exp}} = 5$,
           $\ell_c = 0.5$\,m, $\ell_s^{\mathrm{init}} = [1.0, 1.0]$,
           stratified exploration over $[0, u_M]$.}
  \label{tab:pbe_configs}
  \begin{tabular}{lcccc}
    \toprule
    Config & $z_{\mathrm{rel}}$ & \texttt{base\_position} & $\ell_M$ & $T$ \\
    \midrule
    PBE-B & 1.0\,m & $(0, 0, 0.5)$  & 1.90\,m & 0.85\,s \\
    PBE-C & 1.5\,m & $(0, 0, 1.0)$  & 2.15\,m & 0.95\,s \\
    PBE-D & 2.0\,m & $(0, 0, 1.5)$  & 2.35\,m & 1.00\,s \\
    \bottomrule
  \end{tabular}
\end{table}

\subsection{Implementation}
\label{sec:study3:impl}

\texttt{mc-pilot-pb-elevated/} was created by copying \texttt{mc-pilot-pybullet/}
(which contains all PyBullet infrastructure and interface logic) and adapting it to the elevated hyperparameter
sets.
Noise-specific scripts (\texttt{test\_mc\_pilot\_pb\_B.py},
\texttt{test\_mc\_pilot\_pb\_C.py}) were removed; the zero-noise base
(\texttt{arm\_noise=None}) is used throughout, matching the conditions of Study 1
exactly.
Two new constructor kwargs propagate the pedestal and velocity-limit settings:

\begin{itemize}
  \item \texttt{PyBulletThrowingSystem(base\_position=...,\,vel\_limit\_multiplier=1.5)}
  \item \texttt{ArmController(...,\,vel\_limit\_multiplier=1.5)} via
        \texttt{self.\_qd\_max = \_IIWA\_QD\_MAX~$\times$~vel\_limit\_multiplier}
\end{itemize}

\noindent These changes preserve the default behaviour of \texttt{mc-pilot-pybullet/}
(multiplier $= 1.0$, base at origin) and require no modifications to the RL algorithm
or GP code.

\paragraph{Verification.}
Running \texttt{standalone\_throw.py --direct --speed 3.5 --z\_release 1.0} confirms
$v_{\mathrm{release}} = 3.500$\,m/s and a ball landing at 2.099\,m.
This compares to the NumPy elevated Study 1 verified maximum range of $\approx 1.97$\,m
for Config B---a 5--10\% upward offset attributable to per-step drag application
differences between PyBullet \texttt{applyExternalForce} and the NumPy integrator.
Because the GP is fitted to actual PyBullet throws (not NumPy predictions), it learns
the correct speed-to-range mapping for the PyBullet physics.
The margin $\ell_M = 1.90$\,m $< 2.10$\,m (PyBullet max range) provides 0.2\,m of
coverage room, ensuring all sampled targets remain reachable.

\subsection{PyBullet--NumPy Physics Discrepancy}
\label{sec:study3:discrepancy}

Drag applied via \texttt{applyExternalForce} per simulation step accumulates
small numerical differences compared with the closed-form NumPy midpoint rule.
At $u_M = 3.5$\,m/s and $z_{\mathrm{rel}} = 1.0$\,m the discrepancy is
$\approx 6\%$ in landing range.
This does not impair learning: the GP observes PyBullet ground truth during both
exploration and policy trials, so both the dynamics model and the policy converge
to the PyBullet physics, not the NumPy approximation.
The effect on cross-study comparisons is that trained policies from
\texttt{mc-pilot-elevated/} and \texttt{mc-pilot-pb-elevated/} will produce
slightly different release speeds for the same target, but will achieve comparable
hit rates.

\subsection{Visualisation Pipeline}
\label{sec:study3:viz}

After training completes, \texttt{demo\_pybullet\_gui.py} loads the saved policy
from \texttt{log.pkl} (reading \texttt{parameters\_trial\_list[-1]}, which stores
\{log\_lengthscales, centers, f\_linear.weight\}) and replays $N$ throws in
PyBullet GUI mode.
Each throw renders the full arm trajectory (IK + cubic spline), ball flight with
drag, and landing annotations:
\begin{itemize}
  \item \texttt{addUserDebugLine} trajectory traces (ball path in 3D)
  \item Red sphere at the target location
  \item Green (hit) or orange (miss) sphere at the landing location, with error
        distance printed as GUI text
\end{itemize}
Optional flags: \texttt{--record output.mp4} (via
\texttt{startStateLogging(STATE\_LOGGING\_VIDEO\_MP4)}) and
\texttt{--slow N} to multiply \texttt{time.sleep(dt)} for slower playback.
Command: \texttt{python demo\_pybullet\_gui.py --log\_path results\_mc\_pilot\_pbe\_B/1}.

\subsection{Expected Results}
\label{sec:study3:results}

\needdata{Run Study~3 training scripts for PBE-B/C/D (seed=1, 10 trials each) and
  record: per-trial hit rate, per-throw landing error, convergence cost trajectory.}

Based on the design, PBE-B/C/D are expected to replicate the 10/10 hit rates of
Study 1 Configs B/C/D-Strat: the hyperparameters are identical, the physics
equations are the same (Eq.~35 drag, explicit velocity setting), and the GP learns
from PyBullet ground truth rather than an external model.
The 5--10\% PyBullet drag offset shifts the optimal release speed slightly but does
not affect policy convergence quality.
